Maskinlæring er mye brukt i produktutvikling og produksjon,
men det store antallet tilgjengelige metoder gjør det vanskelig å velge en passende metode,
særlig for brukere med begrenset kompetanse innen maskinlæring.
Relevant kunnskap er spredt over en stor og fragmentert mengde litteratur,
noe som gjør manuelt søk upraktisk.
Denne oppgaven undersøker om ett samlet beslutningsstøttesystem kan anbefale
passende maskinlæringsmetoder og støtte de første stegene i implementeringen.
Tidligere arbeider har en tendens til å behandle disse oppgavene hver for seg.
Tilnærmingen som foreslås her kombinerer ontologibasert kunnskapsrepresentasjon,
semantisk gjenfinning fra vitenskapelig litteratur
og strukturert implementeringsstøtte i ett system.

Et nettbasert system kalt MLguide ble utviklet
innenfor et Design Science Research-rammeverk.
En ontologi over maskinlæringskunnskap leverer et sett med mulige metoder,
der hver metode er koblet til forskningsartiklene som nevner den.
Disse metodene rangeres deretter ved hjelp av semantisk gjenfinning.
Det hentes frem artikler som ligner på brukerens problembeskrivelse,
og en metode rangeres høyt når den nevnes i disse artiklene.
For hvert innsendt problem produseres en rangert liste over anbefalte metoder,
sammen med de høyest rangerte støttende artiklene
og en generert notatbok med startkode for den valgte metoden.

MLguide ble evaluert gjennom brukertesting i to prosjekter i et maskinlæringskurs
og gjennom en eksempelbasert evaluering på konstruerte problemer.
De anbefalte metodene var stort sett relevante for de innsendte problemene,
og dette gjaldt på tvers av maskinlæringsområdene som var dekket.
De støttende artiklene viser grunnlaget for hver anbefaling,
men denne støtten kan være svak,
fordi en metode regnes som støttet så lenge den nevnes i en artikkel,
selv når metoden ikke er artikkelens hovedbidrag.
Flere egenskaper i ontologien har lav dekning,
noe som begrenser ontologiens bidrag.
Å utvide egenskapsdekningen ville gi den en større rolle i anbefalingene.

Resultatene viser at kombinasjonen av disse tre elementene i ett system
kan gi relevante metodeanbefalinger
og støtte de første stegene i implementeringen
for brukere med begrenset maskinlæringskompetanse.
Oppgavens hovedbidrag er selve MLguide,
sammen med koblingen mellom ontologien og forskningsartiklene.